\documentclass[11pt]{article}
\usepackage[margin=1in]{geometry}
\usepackage{amsmath,amssymb}
\usepackage[T1]{fontenc}
\usepackage[utf8]{inputenc}
\usepackage{parskip}
\usepackage{graphicx}
\usepackage{url}
\usepackage{alltt}
\newcommand{\slide}[2]{\subsection*{Slide #1: #2}}
\newcommand{\poll}[2]{\begin{center}\fbox{\parbox{0.85\textwidth}{\centering \textbf{Poll (in-class)} \\ #1 \\ \small #2}}\end{center}}

\title{Amortized Analysis\\[1ex]\large Algorithms COMP.5030 --- structured from NuMarkdown-8B-Thinking OCR (corrected)}
\author{Prof.\ Hugo A.\ Akitaya}
\date{[CLRS ch.\ 17]}

\begin{document}
\maketitle

\section{Introduction}

\slide{1}{Title}
\begin{center}
\textbf{\Large Amortized Analysis}\\[1ex]
Prof.\ Hugo A.\ Akitaya \quad [CLRS ch.\ 17]
\end{center}

\slide{2}{Idea}
\begin{itemize}
  \item \textbf{Na\"ive analysis is not tight:} a simple ``worst-case every time'' estimate can make an algorithm look slower than it really is.
  \item \textbf{Algorithms do many operations:} an algorithm is made of lots of small steps.
  \item \textbf{Some steps are expensive, most are cheap:} occasionally you do a slow/heavy operation, but most of the time you do quick/light ones. If you pretend \emph{every} step is the expensive one, you get an overly pessimistic upper bound.
  \item \textbf{Lazy approach:} don't do big work immediately; postpone it until it's truly needed. This often reduces how often you pay the big cost.
  \item \textbf{Amortization (amortized analysis):} instead of looking at the cost of one operation alone, look at the \textbf{total cost over many operations} and spread it out. This shows the \textbf{average cost per operation} is usually small, even if a few operations are very expensive.
  \item \textbf{Example idea:} resizing a dynamic array is expensive sometimes, but because it happens rarely, the \textbf{average} cost per insert is still small.
\end{itemize}

\section{Example 1: Incrementing a Binary Counter}

\slide{3}{The data structure}
\begin{itemize}
  \item Binary counter $A[1..k]$ with $k$ bits stores a number $x$ with $0 \le x \le 2^k - 1$ ($x = 0$ at the beginning):
        \[ x = \sum_{i=1}^{k} A[i] \cdot 2^{i-1}. \]
  \item Each bit has a different \emph{place value}; the array stores the number \textbf{least significant bit first}:
        \begin{center}
        \begin{tabular}{ccc}
          \hline
          Array position & Bit & Value \\
          \hline
          $A[1]$ & 0 or 1 & $2^0 = 1$ \\
          $A[2]$ & 0 or 1 & $2^1 = 2$ \\
          $A[3]$ & 0 or 1 & $2^2 = 4$ \\
          $A[4]$ & 0 or 1 & $2^3 = 8$ \\
          \hline
        \end{tabular}
        \end{center}
        Example: $A = [1,1,0,1]$ gives $x = 1(2^0) + 1(2^1) + 0(2^2) + 1(2^3) = 1 + 2 + 0 + 8 = 11$. Visually one would write $1011_2$, but the array is $[1,1,0,1]$.
  \item Why $2^k$ values? Each of the $k$ bits has 2 possibilities, so there are $2^k$ combinations; starting at $0$, the largest value is $2^k - 1$. (3 bits: $000, 001, \ldots, 111$ represent $0, 1, \ldots, 7$.)
  \item How do we add one to $x$? Binary addition has a \textbf{carry}, just like $199 + 1 = 200$ in decimal. Start with the least significant bit $A[1]$:
  \begin{itemize}
    \item \textbf{Case 1: $A[1] = 0$.} Then $0 + 1 = 1$: flip it and finish (only \emph{one bit} changes), e.g.\ $[0,1,0,0] \to [1,1,0,0]$.
    \item \textbf{Case 2: $A[1] = 1$.} Then $1 + 1 = 10_2$: set the bit to $0$ and carry the $1$ to the next bit, e.g.\ $[1,0,0,0] \to [0,1,0,0]$.
    \item If the next bit is also $1$, the carry keeps travelling: $[1,1,1,0]$ (value $7$) becomes $[0,0,0,1]$ (value $8$) --- four bits modified.
  \end{itemize}
  \item Pseudocode:
        \begin{alltt}
\textsc{Increment}(A):
    i = 1
    while i <= k and A[i] == 1:
        A[i] = 0
        i = i + 1
    if i <= k:
        A[i] = 1
        \end{alltt}
  \item Why it stops: the first $0$ encountered becomes $1$; the $1$s before it become $0$s; everything after is unchanged: $\underbrace{11\ldots1}_{\text{become 0}}\,\underbrace{0}_{\text{becomes 1}} \to \underbrace{00\ldots0}_{\text{carry}}\,1$.
  \item Two questions: worst-case and best-case cost \emph{per operation}?
\end{itemize}

\slide{4}{Demo}
Binary counter increment demo (interactive animation).

\slide{5}{Worst and best case}
\begin{itemize}
  \item Best case: $A[1] = 0$ already --- only one bit changes, so cost $O(1)$ (essentially one bit update), e.g.\ $[0,1,1,0,1] \to [1,1,1,0,1]$.
  \item Worst case: all $k$ bits are $1$ (e.g.\ $[1,1,1,1,1] = 31$; adding $1$ gives $32 = 100000_2$, which wraps to $00000_2$ with only 5 bits) $\to$ the carry propagates through all $k$ bits, cost $\Theta(k)$. \hfill (slide: $\Omega(k)$)
  \item Odometer intuition (decimal): $0009 \to 0010$ changes two digits, $0010 \to 0011$ changes one, but $0099 \to 0100$ changes three and $0999 \to 1000$ changes four. The binary counter behaves the same way: most operations are cheap, occasionally one is expensive --- \textbf{worst-case cost $\ne$ amortized cost}.
  \item Cost of performing $n$ operations:
  \begin{itemize}
    \item Na\"ive analysis: $O(nk)$.
    \item Intuition: ``good cases'' happen much more often than ``bad cases''. Watching consecutive increments from $0$, the numbers of bits changed are $1,2,1,3,1,2,1,4,\ldots$ --- most increments change only one or two bits.
    \item \textbf{Amortized:} $O(n)$ total cost.
  \end{itemize}
  \item Three answers to remember:
        \begin{center}
        \begin{tabular}{lc}
          \hline
          Analysis & Cost \\
          \hline
          Best case for one increment & $O(1)$ \\
          Worst case for one increment & $O(k)$ \\
          Amortized cost over many increments & $O(1)$ \\
          \hline
        \end{tabular}
        \end{center}
\end{itemize}

\subsection*{Aggregate analysis}

\slide{6}{Aggregate analysis: questions}
\begin{itemize}
  \item Strategy: do \emph{not} analyze each operation in isolation --- instead compute the \textbf{total} cost of a sequence of $n$ operations, then divide by $n$.
  \item The cost of one increment equals the number of bits that flip in it. So ask, for each bit: how many times does it flip over $n$ increments?
  \item Sum up the total cost:
  \begin{itemize}
    \item Bit $A[1]$ changes how many times? (Every operation --- it flips $1 \to 0$ or $0 \to 1$ on \emph{every} increment, so $\approx n$ times.)
    \item How about $A[2]$? \hfill (pair discussion) --- \emph{answer:} it flips every 2nd increment, i.e., whenever $A[1]$ was 1 and carries into it: increments $2, 4, 6, \ldots$, so $\approx n/2$ times.
    \item Following the same pattern: $A[3]$ flips every 4th increment ($\approx n/4$), $A[4]$ every 8th ($\approx n/8$), and in general $A[i]$ every $2^{i-1}$-th increment ($\approx n/2^{i-1}$).
  \end{itemize}
\end{itemize}

\slide{7}{Aggregate analysis: solution}
\begin{itemize}
  \item Why does $A[i]$ flip only then? An increment flips $A[i]$ only when the carry reaches it, i.e., only when $A[1..i-1]$ are all 1s before the increment. Among $n$ consecutive increments starting from $0$, that condition holds for the $i$th bit at most $\lfloor n / 2^{i-1} \rfloor$ times (the counter wraps through its low $i$ bits once every $2^{i-1}$ increments).
  \item Per-bit flip counts over $n$ operations:
        \[ A[1]: n, \quad A[2]: \approx n/2, \quad A[3]: \approx n/4, \quad A[4]: \approx n/8, \quad \ldots \]
        Concretely (8 increments from $0$): bits flipped per increment are $1, 2, 1, 3, 1, 2, 1, 4$ --- a total of $15 = 8 + 4 + 2 + 1$ flips, matching $n + n/2 + n/4 + n/8$.
  \item Total cost (summing over all $k$ bits):
        \[ \sum_{i=1}^{k} \left\lfloor \frac{n}{2^{i-1}} \right\rfloor \le n \sum_{j=0}^{\infty} \frac{1}{2^{j}} = n \cdot \frac{1}{1 - \tfrac12} = 2n \quad \text{(geometric series)}. \]
        Since $1 + \tfrac12 + \tfrac14 + \tfrac18 + \cdots < 2$, the total cost of $n$ increments is $< 2n = O(n)$, \emph{regardless of $k$}.
  \item Amortized cost $=$ total cost $/$ number of operations $= 2n/n = 2 = O(1)$.
  \item Takeaway: even though a \emph{single} increment can cost $\Theta(k)$, the \emph{average} cost over any sequence of increments is $O(1)$ --- the expensive ones are so rare that the geometric series absorbs them.
\end{itemize}

\subsection*{Accounting method}

\slide{8}{Accounting method}
\begin{itemize}
  \item Main idea: \textbf{charge cheap operations a little extra, save that extra amount, and use the savings to pay for expensive operations later.}
  \item For operation $i$: $c_i$ = \textbf{actual cost} (what it really does), $\hat{c}_i$ = \textbf{amortized cost} (what we assign/charge). They don't have to be equal, e.g.\ cheap ops $c_i = 1$ charged $\hat{c}_i = 2$; an expensive op $c_i = 4$ still charged only $\hat{c}_i = 2$.
  \item Bank-account intuition: amortized cost $=$ money you charge; actual cost $=$ money you actually spend. If $\hat{c}_i > c_i$, the leftover becomes \textbf{credit}; later, when $c_i > \hat{c}_i$, accumulated credit covers the difference.
  \item Intuition: leave coins on the data structure to pay later for expensive operations. Each element can have a coin sitting on it --- coins aren't real money, just a bookkeeping device proving we have enough saved credit.
  \item Decide the amortized cost $\hat{c}_i$ of the $i$th operation so that
  \begin{itemize}
    \item $\hat{c}_i \ge c_i$ for cheap operations (saving for a rainy day; the extra $\hat{c}_i - c_i$ becomes credit),
    \item expensive operations may have $\hat{c}_i < c_i$ (shortfall paid from savings),
    \item $\sum_{i=1}^{n} \hat{c}_i \ge \sum_{i=1}^{n} c_i$, i.e.\ $\sum_{i=1}^{n} (\hat{c}_i - c_i) \ge 0$: the total charged must be at least the total actual cost --- we never owe anything in the end (expensive operations are paid from the reserve in the bank). $\hat{c}_i - c_i$ is exactly the change in the bank balance.
  \end{itemize}
  \item Caution: you can't pick any charge you like. Example: actual costs $1, 1, 1, 4, 1, 1, 1, 8$ total $18$, but charging $2$ each gives only $16$ --- not enough. You must \emph{prove} the chosen charge always leaves enough credit.
  \item Warning: amortized analysis $\ne$ average-case analysis. Average-case involves probabilities over inputs; amortized analysis needs \textbf{no randomness} --- it holds for \emph{every} sequence of $n$ operations, even the worst possible one.
\end{itemize}

\slide{9}{Accounting method: binary counter}
\begin{itemize}
  \item Charge $\hat{c}_i = 2$ coins for each increment. Why 2? Every increment (1) flips $t$ ones $\to$ zeros and (2) flips one $0 \to 1$; the newly created 1 can hold a coin.
  \item Invariant: every bit valued 1 has a coin stored on it (so credit $\Phi_i$ $=$ number of saved coins $=$ number of 1s $\ge 0$ at all times --- we never borrow imaginary money from the future).
  \begin{itemize}
    \item $A[1]$ changes $0 \to 1$: pay 1 coin for the change and store 1 coin on $A[1]$.
    \item $A[1]$ changes $1 \to 0$: let $A[j]$ be the bit that changes $0 \to 1$, so $A[1..j-1]$ are all 1s that become 0. Use the $j-1$ stored coins to flip $A[1..j-1]$; from the 2 new coins, pay 1 to change $A[j]$ and store 1 on $A[j]$.
  \end{itemize}
  \item Walkthrough (charging 2 each time): $0000 \to 0001$ ($c_1 = 1$: spend 1, save 1 on the new 1); $0001 \to 0010$ ($c_2 = 2$: 2 charged $+$ the stored coin on the old 1 pay for both flips); $0010 \to 0011$ ($c_3 = 1$: spend 1, save 1); $0011 \to 0100$ ($c_4 = 3 > 2$: the coins stored on the two 1s pay for the carry --- the expensive operation was \textbf{prepaid by earlier cheap ones}).
  \item General increment with $t$ consecutive 1s: actual cost $c_i = t + 1$ (can be $O(k)$), but the $t$ stored coins pay for the $t$ bits flipping $1 \to 0$, and the 2 new coins pay for the new 1 (1 to spend, 1 to save). All work is paid for, so $\hat{c}_i = 2 = O(1)$ safely.
  \item Every $1 \to 0$ flip is paid by a coin stored when that bit earlier went $0 \to 1$ --- expensive carries never come out of nowhere.
  \item The bank is never negative (the number of 1s is nonnegative), hence
        \[ \sum_{i=1}^{n} \hat{c}_i \;\ge\; \sum_{i=1}^{n} c_i. \]
  \item Big picture: actual cost $c_i = t + 1 = O(k)$; amortized cost $\hat{c}_i = 2 = O(1)$; total over $n$ increments $\sum c_i = O(n)$, so the average is $O(n)/n = O(1)$.
  \item Exam checklist: (1) give each op an amortized charge $\hat{c}_i$; (2) cheap ops overcharged; (3) the difference becomes credit; (4) expensive ops use the credit; (5) credit never negative; (6) show $\sum \hat{c}_i \ge \sum c_i$; (7) for the binary counter, $\hat{c}_i = 2$ gives $O(1)$ amortized even though one increment can cost $O(k)$.
\end{itemize}

\subsection*{Potential method}

\slide{10}{Potential method}
\begin{itemize}
  \item The data structure $D$ changes with each operation: $D_0$ is the initial state, $D_i$ the state after the $i$th operation.
  \item Define a potential function $\Phi : D \to \mathbb{R}$.
  \item Common (not required) conditions: $\Phi(D_0) = 0$ and $\Phi(D_i) \ge 0$ for all $i$.
  \item Required: $\Phi(D_n) \ge \Phi(D_0)$.
  \item Amortized cost of the $i$th operation:
        \[ \hat{c}_i = c_i + \Delta\Phi = c_i + \Phi(D_i) - \Phi(D_{i-1}). \]
\end{itemize}

\slide{11}{Potential method: binary counter}
\begin{itemize}
  \item Define $\Phi(D)$ $=$ number of 1s in $D$ (number of coins in the bank); $\Phi(D_0) = 0$, $\Phi(D_i) \ge 0$.
  \item Let $t_i$ be the number of 1s turned to 0 by the $i$th increment. The real cost is $c_i = t_i + 1$, and the potential change is $\Delta\Phi = 1 - t_i$:
        \[ \hat{c}_i = c_i + \Delta\Phi = (t_i + 1) + (1 - t_i) = 2. \]
\end{itemize}

\slide{12}{Demo}
Binary counter demo continued (interactive animation).

\section{Example 2: Dynamic Array --- Push Only}

\slide{13}{Dynamic array (push only)}
\begin{itemize}
  \item A \textbf{dynamic array} is an array that can \textbf{grow} when you keep adding (``push'') items, like a \textbf{stack}.
  \item Goal: implement a stack with an array.
  \item Problem with normal arrays: you must choose the size early. Arrays are classically allocated statically (maybe before we know the size of the data):
  \begin{itemize}
    \item Allocate too much space $\to$ wasted space (wasted memory).
    \item Allocate too little space $\to$ wasted time: you must create a bigger array and copy everything into it.
  \end{itemize}
  \item Idea: keep the array \textbf{about half full} --- maintain the load factor at a constant (around $1/2$?). That way you don't waste too much space, and you don't resize too often.
  \item Operations: \textsc{Push} --- worst case? \hfill (pair discussion) --- usually push is \textbf{fast} ($O(1)$: just place the element); worst case, when the array is full, push is \textbf{slow} because of copying ($O(num)$). But doubling makes this rare, so the average push is still fast (amortized $O(1)$).
\end{itemize}

\slide{14}{Insert pseudocode}
\begin{itemize}
  \item What the variables mean: $A$ = the underlying array (storage); $size$ = the \textbf{capacity} of $A$ (how many slots exist), initially 0; $num$ = the \textbf{number of elements currently stored} (how many slots are used), initially 0.
  \item $\textsc{Push}(x)$ in simple steps:
  \begin{itemize}
    \item If this is the first insert ($size = 0$, edge case): create an array (often size $= 1$).
    \item If the array is \textbf{full} ($num = size$): make a \textbf{new array} with \textbf{double} capacity $2 \cdot size$, \textbf{copy} all old elements into it, and replace $A$ with the new array.
    \item Put the new element in the next free spot: $A[num] = x$.
    \item Increase $num$ by 1 (and update $size$ if a resize happened).
  \end{itemize}
  \item Cost: cheap push $O(1)$; push with resize $O(num)$ --- but resizes happen rarely, so amortized $O(1)$.
\end{itemize}

\slide{15}{Aggregate method}
\begin{itemize}
  \item Main idea: most insertions are cheap, occasionally one is expensive, but averaged across many insertions each is still $O(1)$.
  \item What happens when we push? Starting from capacity 1, the array doubles when full:
        \begin{alltt}
capacity 1: [A]
push B -> resize: [A, B]                        (capacity 2)
push C -> resize: [A, B, C, \_]                 (capacity 4)
push D -> no resize: [A, B, C, D]
push E -> resize: [A, B, C, D, E, \_, \_, \_]   (capacity 8)
        \end{alltt}
  \item Perform $n$ inserts. Two kinds of pushes:
  \begin{itemize}
    \item \textbf{Cheap push} (free space available): just add the item at the tail, e.g.\ $[A,B,\_,\_] \to [A,B,C,\_]$; cost $O(1)$ --- happens often.
    \item \textbf{Expensive push} (array full): create a larger array and copy everything, e.g.\ copy $[A,B,C,D]$ into a new array of capacity 8; cost $O(i)$ because there may be $i$ elements to copy --- happens only if $i = 2^k$ for some integer $k$ (resizes occur at sizes $1, 2, 4, 8, \ldots$).
  \end{itemize}
  \item Aggregate method: don't worry about the cost of one push --- add up the cost of all $n$ pushes. Most pushes cost 1, contributing $O(n)$; the total copying work is $1 + 2 + 4 + 8 + \cdots$ (e.g.\ for 16 pushes: $1 + 2 + 4 + 8 = 15 < 16$).
  \item Total cost:
        \[ O(n) + \sum_{k=0}^{\lfloor \log_2 n \rfloor} 2^k \;\le\; O(n) + 2n \;=\; O(n). \]
        In words: if you add all the copying caused by resizing, the total copying is at most roughly $2n$.
  \item Average cost per push $= O(n)/n = O(1)$.
\end{itemize}

\slide{16}{Accounting method}
\begin{itemize}
  \item Goal: show that even though resizing is expensive, the \emph{average} cost per push is $O(1)$ using the ``coin'' (accounting) idea.
  \item Perform $n$ pushes; charge $\hat{c}_i = 3$ per push:
  \begin{itemize}
    \item \textbf{1 coin} pays for the actual insertion (write the element at the end).
    \item The remaining \textbf{2 coins} are \textbf{saved} to pay for future copying during a resize.
  \end{itemize}
  \item Two equivalent ways to ``store'' the saved coins (the difference is only \emph{where} the coins sit; every push still uses 1 coin for insertion):
  \begin{itemize}
    \item \textbf{Version 1 (simplest):} save 2 coins on \emph{every} element. After 4 pushes (capacity 4): $4 \times 2 = 8$ saved coins.
    \item \textbf{Version 2 (like the slide):} save 2 coins only on elements in the \emph{last half} of the array. For capacity 4, last half $=$ positions 3 and 4 $\to 2 \times 2 = 4$ saved coins (savings concentrated, not spread).
  \end{itemize}
  \item Invariant (version 2): every element in the last half of the array stores two coins; the bank is always positive --- we never spend more coins than we have saved, so we never go into ``coin debt.'' Example (capacity 8, after resizing from 4): the next four cheap pushes E, F, G, H land in the last half, each storing 2 coins $\to 4 \times 2 = 8$ saved coins, exactly enough to copy all 8 elements at the next resize. In general, capacity $n$: last half has $\tfrac n2$ elements $\times 2$ coins $= n$ coins for copying the $n$ existing elements.
  \item \textbf{Cheap insertion:} use 1 coin to pay for the push and deposit 2 coins on the new position.
  \item \textbf{Expensive insertion:} by the invariant there are $\tfrac{n}{2} \cdot 2 = n$ coins stored; use them to pay for the array expansion. We still have the 3 new coins: a cheap push into the new array.
  \item What happens on the 5th push (array full, resize $4 \to 8$):
  \begin{enumerate}
    \item \textbf{Resize + copy:} create a new array of size 8 and copy the 4 elements --- copying costs 4 operations, paid using the saved coins.
    \item \textbf{Insert 5th element:} put it into the new array (cost 1 coin).
    \item \textbf{Save again:} from the 3 coins charged, after paying 1, save 2 coins (toward the next resize).
  \end{enumerate}
  \item Key idea: resizes are paid for by coins saved from many earlier cheap pushes, so the average cost stays constant.
\end{itemize}

\slide{17}{Potential method}
\begin{itemize}
  \item Both this and the accounting method prove the same result: pushing into a dynamic array has \textbf{amortized cost $O(1)$}, even though resizing is occasionally expensive.
  \item The potential method is basically the same savings idea, but instead of imaginary coins, the amount saved is represented by a mathematical function, the \textbf{potential} $\Phi$. Mentally: $\Phi = $ \emph{saved work/coins for the future}.
  \item Define $\Phi$ with $\Phi(A_0) = 0$.
  \item Idea: before an expensive push we have just enough coins, i.e., $i$ coins for the $i$th push; after a resize the potential is (nearly) depleted. Let $size_i$ denote the array size after the $i$th operation ($i$ = number of elements after $i$ pushes):
        \[ \Phi(A_i) = \left( i - \frac{size_i}{2} \right) \cdot 2 = 2i - size_i. \]
  \item Example: resized to capacity 8 with 4 elements $[A,B,C,D,\_,\_,\_,\_]$: $\Phi = 2(4) - 8 = 0$ --- immediately after a resize the potential is $0$.
  \item Insert $E$: $\Phi = 2(5) - 8 = 2$. A cheap push saved exactly 2 coins in the accounting method, and here the potential increases by 2 --- \emph{the same idea}! Insert $F$: $\Phi = 2(6) - 8 = 4$.
  \item Growth between resizes (capacity 8):
        \begin{center}
        \begin{tabular}{ccc}
          \hline
          Elements $i$ & Capacity $size_i$ & Potential $2i - size_i$ \\
          \hline
          4 & 8 & 0 \\
          5 & 8 & 2 \\
          6 & 8 & 4 \\
          7 & 8 & 6 \\
          8 & 8 & 8 \\
          \hline
        \end{tabular}
        \end{center}
        Just before the next resize ($[A,\ldots,H]$ full), $\Phi = 8$: enough accumulated ``saved work'' to pay for moving the 8 existing elements.
  \item Then the expensive push happens: the array is full, capacity doubles $8 \to 16$, the accumulated potential accounts for the copying, and after the resize the potential resets ($\approx 0$) and saving begins again.
  \item Accounting vs.\ potential: the accounting method thinks in \emph{coins} (charge 3, spend 1, save 2, use savings for the resize); the potential method thinks in \emph{stored energy} (potential increases with cheap pushes, is consumed by the resize). Saved coins $\approx$ potential.
  \item Why do this? Looking at one resize alone suggests push is $O(n)$, but that's misleading: resizing doesn't happen on every push. By saving a little credit during cheap operations we pay for the rare expensive ones, so the amortized cost per push is $O(1)$.
  \item Amortized costs: \hfill (pair discussion)
\end{itemize}

\slide{18}{Poll: potential of a given array}
\poll{What is the potential of the following array?}{Forms: \url{https://forms.office.com/r/AZbBVrvYEj}}

\section{Dynamic Array --- Push and Pop}

\slide{19}{Dynamic array (push + pop)}
\begin{itemize}
  \item Main trick: \textbf{grow when full, but don't shrink too quickly when deleting.}
  \item Reminder: $size$ = capacity of the array; $num$ = number of elements currently stored (e.g.\ $[A,B,C,\_,\_,\_,\_,\_]$: $num = 3$, $size = 8$).
  \item \textbf{Load factor} $=$ how full the array is: $\alpha = num/size$. Example $[A,B,C,D,\_,\_,\_,\_]$: $\tfrac48 = \tfrac12 = 50\%$ full.
  \item Goal: maintain the load factor at least a constant (around how much?).
  \item Operations:
  \begin{itemize}
    \item \textsc{Push}: same as before --- cheap $O(1)$ if there is space; if full, double the array and copy (expensive, but rare).
    \item \textsc{Pop}: removes the last element, e.g.\ $[A,B,C,D] \to [A,B,C,\_]$; normally very cheap ($O(1)$). Worst case?
  \end{itemize}
  \item New problem: capacity 16 but only 2 elements left after many pops $\to$ a large array holding almost nothing = wasted memory, so at some point we should shrink.
  \item Pair: what happens if you downsize by half whenever the load factor drops below $1/2$? (Threatening: thrashing at the boundary.) Popping drops it slightly below $1/2 \to$ shrink; pushing fills it up $\to$ double; pop $\to$ shrink; push $\to$ grow; \ldots every operation pays for a resize that copies many elements.
  \item Solution: deliberately leave a \textbf{gap} between the two resize conditions --- grow when 100\% full, shrink only below 25\% full:
        \begin{center}
        \begin{tabular}{ccc}
          0\% & 25\% & 50\% \quad\quad 100\% \\
          \hline
          \multicolumn{2}{c}{\textsc{shrink} (below $1/4$)} & \textsc{grow} (full) \\
        \end{tabular}
        \end{center}
        Bus analogy: when the bus is completely full, get one twice as big; but if a few passengers leave, you don't immediately switch to a smaller bus --- you wait until it's very empty (around $1/4$ full).
  \item Now $num_i \ne i$ in general: $num_i$ denotes the number of elements after the $i$th operation.
\end{itemize}

\slide{20}{Delete pseudocode}
\begin{itemize}
  \item State: array $A$ (initially \textsc{nil}), integer $size$ (capacity, initially 0), integer $num$ (number of elements, initially 0).
  \item $\textsc{Pop}()$ step by step:
  \begin{itemize}
    \item \textbf{Edge case when $num = 1$}: removing it makes the array empty, so special handling is needed.
    \item \textbf{Check if the array is too empty}: $num - 1 < \tfrac14 \cdot size$, where $num - 1$ = ``how many elements will remain after \textsc{Pop}''. If that would be less than $25\%$ of capacity, shrink.
    \item \textbf{Shrink}: create a new array of length $size/2$ and copy the data (e.g.\ capacity $16 \to 8$; with only 3 elements in a capacity-16 array, $\tfrac14 \cdot 16 = 4$, and $3 < 4$, so we shrink to $[A,B,C,\_,\_,\_,\_,\_]$ --- much less wasted memory).
    \item \textbf{Delete the last element}: $A[num] = \textsc{nil}$ clears the position being removed.
    \item \textbf{Update $size$ and $num$}: $num$ decreases (one fewer element); $size$ decreases if the array was resized.
  \end{itemize}
  \item Exam summary: \textsc{Push} --- if full, \textbf{double} the array; \textsc{Pop} --- if less than $1/4$ full, \textbf{cut size in half}; the reason for $1/4$ instead of $1/2$ is \textbf{to avoid repeatedly growing and shrinking}. $\tfrac{num}{size}$ = ``how full is my array?''
\end{itemize}

\slide{21}{Potential method (push + pop)}
\begin{itemize}
  \item Basic idea: \textbf{potential $=$ saved credit for future expensive resizing}. Cheap operations build up enough credit to pay for the occasional expensive resize.
  \item Terms: $num$ = number of elements actually inside the array; $size$ = total capacity; $\alpha = num/size$ = load factor (how full the array is, e.g.\ $num = 4$, $size = 8 \Rightarrow \alpha = 1/2$).
  \item Two expensive situations to prepare for, i.e., potential must work in \emph{both directions}:
  \begin{itemize}
    \item \textbf{Expensive push:} array full $\to$ double and copy (need credit when too full).
    \item \textbf{Expensive pop:} array below $1/4$ full $\to$ halve and copy (need credit when too empty).
  \end{itemize}
  \noindent So: too empty $\to$ shrink (credit needed) \quad $\bullet$ \quad around $1/2$ full $\to$ no resize (comfortable) \quad $\bullet$ \quad too full $\to$ grow (credit needed).
  \item Define $\Phi$ with $\Phi(A_0) = 0$; after a resize operation the potential is $0$.
  \item Build up potential so that the next expensive operation is prepaid:
  \begin{itemize}
    \item if the next operation doubles $A$ (i.e., $A_i$ is full), the potential reaches $size_i$;
    \item if the next operation halves $A$ (i.e., $A_i$ is $\tfrac{1}{4}$ full), the potential reaches $size_i/4$.
  \end{itemize}
  \item With load factor $\alpha_i = num_i / size_i$ (one formula when at least half full, another when less than half full):
        \[
          \Phi(A_i) =
          \begin{cases}
            2 \cdot num_i - size_i & \text{if } \alpha_i \ge \tfrac{1}{2},\\[0.5ex]
            size_i/2 - num_i & \text{if } \alpha_i < \tfrac{1}{2}.
          \end{cases}
        \]
  \item \textbf{Growing direction} ($\alpha \ge 1/2$, $\Phi = 2\,num - size$): with $size = 8$, potential is $0, 2, 4, 6, 8$ for $num = 4, 5, 6, 7, 8$ --- it rises as the array approaches full ($\Phi = 8$ when full, ready to pay for copying 8 elements).
  \item After doubling: $num = 8$, $size = 16$, exactly half full $\to \Phi = 2(8) - 16 = 0$ --- saved credit was used for the resize and savings are back to zero.
  \item \textbf{Shrinking direction} ($\alpha < 1/2$, $\Phi = size/2 - num$): with $size = 16$, potential is $0, 1, 2, 3, 4$ for $num = 8, 7, 6, 5, 4$ --- the emptier the array, the more credit saved to pay for shrinking (near $1/4$ full, $\Phi = 4 = size/4$).
  \item Battery picture: potential is smallest around 50\% full and grows as we approach a resize (toward 100\%: save for growing; toward 25\%: save for shrinking). Think of it as a savings account: pushes/pops deposit, a resize withdraws.
  \item Exam cheat sheet: $num$ = elements, $size$ = capacity, $\alpha = num/size$ = how full, $\Phi$ = saved credit; near 100\% $\to$ save for growing; near 25\% $\to$ save for shrinking; after a resize credit drops back down; result: push and pop are $O(1)$ amortized.
\end{itemize}

\slide{22}{Amortized cost: case analysis}
We must show that $\hat{c}_i \in O(1)$ in each case --- i.e., individual operations can cost $O(n)$ (copying during resize), but resizes don't happen often, so the \emph{amortized} cost of each push or pop is constant:
\begin{itemize}
  \item \textbf{Case 1, cheap push:} there is room $\to$ just insert, no resize (analyze for $\alpha_i \ge 1/2$ and $\alpha_i < 1/2$).
  \item \textbf{Case 2, expensive push:} array full $\to$ double, copy, insert --- saved potential pays for the copying.
  \item \textbf{Case 3, cheap pop:} remove last element, no resize (analyze for $\alpha_i \ge 1/2$ and $\alpha_i < 1/2$).
  \item \textbf{Case 4, expensive pop:} array too empty $\to$ shrink, copy --- saved potential again pays.
\end{itemize}
(The details were worked out on the board; the two expensive cases are derived on Slides 23--24.)

\slide{23}{Amortized cost: expensive push}
\textbf{Push with resize (expensive):}
\begin{itemize}
  \item The array was full: \quad $num_{i-1} = size_{i-1}$ \hfill (1)
  \item We double the array: \quad $size_i = 2 \cdot size_{i-1}$ \hfill (2)
  \item Real cost of copying the array and adding one element: \quad $c_i = num_{i-1} + 1$ \hfill (3)
  \item Potential before:
        \[ \Phi_{i-1} = 2\,num_{i-1} - size_{i-1} = num_{i-1} \quad \text{(using (1))} \]
  \item Potential after:
        \[ \Phi_i = 2\,num_i - size_i = 2 \quad \text{(using (1), (2), and } num_i = num_{i-1} + 1\text{)} \]
  \item Amortized cost:
        \[ \hat{c}_i = c_i + \Delta\Phi = num_{i-1} + 1 + (2 - num_{i-1}) = 3. \]
\end{itemize}

\slide{24}{Amortized cost: expensive pop}
\textbf{Pop with resize (expensive):}
\begin{itemize}
  \item The array was $\tfrac{1}{4}$ full: \quad $4\,num_{i-1} = size_{i-1}$ \hfill (1)
  \item We halve the array: \quad $size_{i-1} = 2 \cdot size_i$ \hfill (2)
  \item Real cost of copying the array and popping one element: \quad $c_i = num_{i-1} + 1$ \hfill (3)
  \item Potential before:
        \[ \Phi_{i-1} = size_{i-1}/2 - num_{i-1} = num_{i-1} \quad \text{(using (1))} \]
  \item Potential after:
        \[ \Phi_i = size_i/2 - num_i = 1 \quad \text{(using (1), (2), and } num_i = num_{i-1} - 1\text{)} \]
  \item Amortized cost:
        \[ \hat{c}_i = c_i + \Delta\Phi = num_{i-1} + 1 + (1 - num_{i-1}) = 2. \]
\end{itemize}

\slide{25}{Poll: potential of a given array}
\poll{What is the potential of the following array? (options: 0, 1, 2, 4)}{Forms: \url{https://forms.office.com/r/e36xVbSVjm}}

\end{document}
